
\documentclass{article}
\usepackage{colortbl}
\usepackage{makecell}
\usepackage{multirow}
\usepackage{supertabular}

\begin{document}

\newcounter{utterance}

\centering \large Interaction Transcript for game `function\_detective', experiment `list\_passive\_labeled\_pairs', episode 75 with mistral{-}large{-}3{-}azure.
\vspace{24pt}

{ \footnotesize  \setcounter{utterance}{1}
\setlength{\tabcolsep}{0pt}
\begin{supertabular}{c@{$\;$}|p{.15\linewidth}@{}p{.15\linewidth}p{.15\linewidth}p{.15\linewidth}p{.15\linewidth}p{.15\linewidth}}
    \# & \multicolumn{2}{c}{Player} && \multicolumn{2}{c}{Game Master} \\
    \hline

    \theutterance \stepcounter{utterance}  
    & & & \multicolumn{4}{p{0.6\linewidth}}{
        \cellcolor[rgb]{0.9,0.9,0.9}{
            \makecell[{{p{\linewidth}}}]{
                \texttt{\tiny{[P1$\langle$GM]}}
                \texttt{I have a hidden Python function for you to figure out.} \\
\\ 
\texttt{Your job is to infer what it does and then submit a correct Python implementation.} \\
\\ 
\texttt{Function signature:} \\
\texttt{(items: List{[}int{]}) {-}> int} \\
\\ 
\texttt{You have 15 turns total.} \\
\\ 
\texttt{I reveal candidate input{-}output pairs one at a time and tell you whether each one is correct.} \\
\\ 
\texttt{You must make your best SOLVE attempt no later than the second{-}to{-}last turn.} \\
\texttt{Do not wait until the very end.} \\
\\ 
\texttt{You can use these actions:} \\
\texttt{{-} JSON format described at the end of the prompt} \\
\texttt{{-} SOLVE: ```python ... ```} \\
\\ 
\texttt{Do not add import statements in your guess.} \\
\texttt{Anything you might need is already available, so just write the function body in hypothesis and return the answer.} \\
\\ 
\texttt{Preloaded examples:} \\
\texttt{1. (({[}5, {-}5, 5{]}, 2), True)} \\
\\ 
\texttt{Scoring:} \\
\texttt{{-} correctness is required for any non{-}zero score} \\
\texttt{{-} in interactive modes, solving with fewer information requests is better} \\
\\ 
\texttt{Input formatting reminders:} \\
\texttt{{-} strings must be quoted} \\
\texttt{{-} booleans must be True/False} \\
\texttt{{-} use commas between multiple values} \\
\texttt{{-} do not write any extra action tags in the same message} \\
\\ 
\texttt{Format example:} \\
\texttt{You write:} \\
\texttt{\{'current\_mode': <probing/confirming/don't know; choose whichever best describes your current state in the task>, 'input\_rationale': <brief description of why you chose this hypothesis>, 'hypothesis': <either the exact string "unknown" if you have no guess yet OR your current best guess as runnable Python matching the given signature, with a one{-}line docstring stating what it does. Never write placeholder code, if you cannot commit to a guess, use "unknown". Match the given signature.} \\
\texttt{e.g.} \\
\texttt{def f(x):} \\
\texttt{	"""Returns the absolute value of x."""} \\
\texttt{	return x if x >= 0 else {-}x} \\
\\ 
\texttt{Use "unknown" if you have no guess.>, 'confidence\_score': <score between 1{-}100. 100 means highest confidence in your hypothesis>\}} \\
\\ 
\texttt{Output only one valid json per turn.} \\
            }
        }
    }
     \\ \\

    \theutterance \stepcounter{utterance}  
    & \multicolumn{4}{p{0.6\linewidth}}{
        \cellcolor[rgb]{0.561,0.737,0.561}{
            \makecell[{{p{\linewidth}}}]{
                \texttt{\tiny{[P1$\rangle$GM]}}
                \texttt{```json} \\
\texttt{\{} \\
\texttt{  "current\_mode": "probing",} \\
\texttt{  "input\_rationale": "The preloaded example ({[}5, {-}5, 5{]}, 2) suggests the function might be counting the number of positive integers, but the presence of a negative number complicates this. Alternatively, it could be counting the number of times the sign changes or some other pattern. I will test a simple hypothesis first: counting the number of positive integers in the list.",} \\
\texttt{  "hypothesis": "def f(items):\textbackslash{}n    \textbackslash{}"\textbackslash{}"\textbackslash{}"Returns the count of positive integers in the list.\textbackslash{}"\textbackslash{}"\textbackslash{}"\textbackslash{}n    return sum(1 for x in items if x > 0)",} \\
\texttt{  "confidence\_score": 60} \\
\texttt{\}} \\
\texttt{```} \\
            }
        }
    }
    & & \\ \\

    \theutterance \stepcounter{utterance}  
    & & & \multicolumn{4}{p{0.6\linewidth}}{
        \cellcolor[rgb]{0.9,0.9,0.9}{
            \makecell[{{p{\linewidth}}}]{
                \texttt{\tiny{[P1$\langle$GM]}}
                \texttt{OUTPUT: (({[}1, 2, 3{]}, {-}1), False)} \\
            }
        }
    }
     \\ \\

    \theutterance \stepcounter{utterance}  
    & \multicolumn{4}{p{0.6\linewidth}}{
        \cellcolor[rgb]{0.561,0.737,0.561}{
            \makecell[{{p{\linewidth}}}]{
                \texttt{\tiny{[P1$\rangle$GM]}}
                \texttt{```json} \\
\texttt{\{} \\
\texttt{  "current\_mode": "probing",} \\
\texttt{  "input\_rationale": "The previous hypothesis (count of positive integers) is incorrect because the output for {[}1, 2, 3{]} is {-}1, not 3. This suggests the function is not simply counting positives. The preloaded example ({[}5, {-}5, 5{]}, 2) and this new example ({[}1, 2, 3{]}, {-}1) indicate the function might involve some form of alternating sign or cumulative sum behavior. I will test a hypothesis involving the count of sign changes in the list.",} \\
\texttt{  "hypothesis": "def f(items):\textbackslash{}n    \textbackslash{}"\textbackslash{}"\textbackslash{}"Returns the number of sign changes in the list (ignoring zeros).\textbackslash{}"\textbackslash{}"\textbackslash{}"\textbackslash{}n    if not items:\textbackslash{}n        return 0\textbackslash{}n    sign\_changes = 0\textbackslash{}n    prev\_sign = 1 if items{[}0{]} > 0 else {-}1 if items{[}0{]} < 0 else 0\textbackslash{}n    for x in items{[}1:{]}:\textbackslash{}n        current\_sign = 1 if x > 0 else {-}1 if x < 0 else 0\textbackslash{}n        if current\_sign != 0 and prev\_sign != 0 and current\_sign != prev\_sign:\textbackslash{}n            sign\_changes += 1\textbackslash{}n        prev\_sign = current\_sign if current\_sign != 0 else prev\_sign\textbackslash{}n    return sign\_changes",} \\
\texttt{  "confidence\_score": 7} \\
            }
        }
    }
    & & \\ \\

    \theutterance \stepcounter{utterance}  
    & & & \multicolumn{4}{p{0.6\linewidth}}{
        \cellcolor[rgb]{0.9,0.9,0.9}{
            \makecell[{{p{\linewidth}}}]{
                \texttt{\tiny{[P1$\langle$GM]}}
                \texttt{Game Violation: ParseError: Response is not valid JSON.} \\
\texttt{Output ONLY:} \\
\texttt{One JSON with confidence\_score, current\_mode, input\_rationale, and hypothesis keys ...} \\
\texttt{OR} \\
\texttt{SOLVE: ```python ... ```} \\
            }
        }
    }
     \\ \\

    \theutterance \stepcounter{utterance}  
    & \multicolumn{4}{p{0.6\linewidth}}{
        \cellcolor[rgb]{0.561,0.737,0.561}{
            \makecell[{{p{\linewidth}}}]{
                \texttt{\tiny{[P1$\rangle$GM]}}
                \texttt{```json} \\
\texttt{\{} \\
\texttt{  "current\_mode": "probing",} \\
\texttt{  "input\_rationale": "The previous hypothesis about sign changes doesn't fit the preloaded example ({[}5, {-}5, 5{]}, 2), which would have 2 sign changes (5 to {-}5 and {-}5 to 5), but the new example ({[}1, 2, 3{]}, {-}1) contradicts this. I will now test a hypothesis involving the maximum absolute value in the list, but this also doesn't fit the preloaded example. Instead, I suspect the function might be related to the count of local maxima or minima, or perhaps the number of times the sequence alternates between increasing and decreasing.",} \\
\texttt{  "hypothesis": "def f(items):\textbackslash{}n    \textbackslash{}"\textbackslash{}"\textbackslash{}"Returns the number of times the sequence changes from increasing to decreasing or vice versa.\textbackslash{}"\textbackslash{}"\textbackslash{}"\textbackslash{}n    if len(items) < 2:\textbackslash{}n        return 0\textbackslash{}n    changes = 0\textbackslash{}n    for i in range(1, len(items) {-} 1):\textbackslash{}n        prev\_diff = items{[}i{]} {-} items{[}i{-}1{]}\textbackslash{}n        next\_diff = items{[}i+1{]} {-} items{[}i{]}\textbackslash{}n        if (prev\_diff > 0 and next\_diff < 0) or (prev\_diff < 0 and next\_diff > 0):\textbackslash{}n            changes += 1\textbackslash{}n    return changes",} \\
\texttt{  "confidence\_score": 50} \\
\texttt{\}} \\
\texttt{```} \\
            }
        }
    }
    & & \\ \\

    \theutterance \stepcounter{utterance}  
    & & & \multicolumn{4}{p{0.6\linewidth}}{
        \cellcolor[rgb]{0.9,0.9,0.9}{
            \makecell[{{p{\linewidth}}}]{
                \texttt{\tiny{[P1$\langle$GM]}}
                \texttt{OUTPUT: (({[}0{]}, 0), True)} \\
            }
        }
    }
     \\ \\

    \theutterance \stepcounter{utterance}  
    & \multicolumn{4}{p{0.6\linewidth}}{
        \cellcolor[rgb]{0.561,0.737,0.561}{
            \makecell[{{p{\linewidth}}}]{
                \texttt{\tiny{[P1$\rangle$GM]}}
                \texttt{```json} \\
\texttt{\{} \\
\texttt{  "current\_mode": "probing",} \\
\texttt{  "input\_rationale": "The new example ({[}0{]}, 0) suggests the function handles edge cases like empty or single{-}element lists gracefully. The previous hypotheses (sign changes, local maxima/minima) don't align well with all examples. I will now test a hypothesis that the function counts the number of *peaks* in the list, where a peak is an element greater than its neighbors. This fits the preloaded example ({[}5, {-}5, 5{]}, 2) if we consider the first and last elements as peaks (since they are greater than their single neighbor), and it also fits ({[}0{]}, 0) and ({[}1, 2, 3{]}, {-}1) if we assume {-}1 is a sentinel for 'no peaks'. However, this is still uncertain.",} \\
\texttt{  "hypothesis": "def f(items):\textbackslash{}n    \textbackslash{}"\textbackslash{}"\textbackslash{}"Returns the number of peaks in the list, where a peak is an element greater than its neighbors.\textbackslash{}"\textbackslash{}"\textbackslash{}"\textbackslash{}n    if len(items) == 1:\textbackslash{}n        return 0\textbackslash{}n    peaks = 0\textbackslash{}n    for i in range(len(items)):\textbackslash{}n        if i == 0:\textbackslash{}n            if items{[}i{]} > items{[}i+1{]}:\textbackslash{}n                peaks += 1\textbackslash{}n        elif i == len(items) {-} 1:\textbackslash{}n            if items{[}i{]} > items{[}i{-}1} \\
            }
        }
    }
    & & \\ \\

    \theutterance \stepcounter{utterance}  
    & & & \multicolumn{4}{p{0.6\linewidth}}{
        \cellcolor[rgb]{0.9,0.9,0.9}{
            \makecell[{{p{\linewidth}}}]{
                \texttt{\tiny{[P1$\langle$GM]}}
                \texttt{Game Violation: ParseError: Response is not valid JSON.} \\
\texttt{Output ONLY:} \\
\texttt{One JSON with confidence\_score, current\_mode, input\_rationale, and hypothesis keys ...} \\
\texttt{OR} \\
\texttt{SOLVE: ```python ... ```} \\
            }
        }
    }
     \\ \\

    \theutterance \stepcounter{utterance}  
    & \multicolumn{4}{p{0.6\linewidth}}{
        \cellcolor[rgb]{0.561,0.737,0.561}{
            \makecell[{{p{\linewidth}}}]{
                \texttt{\tiny{[P1$\rangle$GM]}}
                \texttt{```json} \\
\texttt{\{} \\
\texttt{  "current\_mode": "probing",} \\
\texttt{  "input\_rationale": "The previous hypotheses don't fully explain all examples. The preloaded example ({[}5, {-}5, 5{]}, 2) suggests the function counts the number of times the sequence reaches a 'peak' or 'valley' (local maxima/minima). The new example ({[}0{]}, 0) fits this if we return 0 for single{-}element lists. The example ({[}1, 2, 3{]}, {-}1) contradicts this, so perhaps the function returns {-}1 for strictly increasing/decreasing sequences. I will now test a refined hypothesis: count the number of local maxima (peaks) and minima (valleys), but return {-}1 if the sequence is strictly increasing or decreasing.",} \\
\texttt{  "hypothesis": "def f(items):\textbackslash{}n    \textbackslash{}"\textbackslash{}"\textbackslash{}"Returns the number of local maxima and minima in the list, or {-}1 if the sequence is strictly increasing or decreasing.\textbackslash{}"\textbackslash{}"\textbackslash{}"\textbackslash{}n    if len(items) <= 1:\textbackslash{}n        return 0\textbackslash{}n    is\_increasing = True\textbackslash{}n    is\_decreasing = True\textbackslash{}n    extrema = 0\textbackslash{}n    for i in range(1, len(items)):\textbackslash{}n        if items{[}i{]} > items{[}i{-}1{]}:\textbackslash{}n            is\_decreasing = False\textbackslash{}n        elif items{[}i{]} < items{[}i{-}1{]}:\textbackslash{}n            is\_increasing = False\textbackslash{}n    if is} \\
            }
        }
    }
    & & \\ \\

    \theutterance \stepcounter{utterance}  
    & & & \multicolumn{4}{p{0.6\linewidth}}{
        \cellcolor[rgb]{0.9,0.9,0.9}{
            \makecell[{{p{\linewidth}}}]{
                \texttt{\tiny{[P1$\langle$GM]}}
                \texttt{Game Violation: ParseError: Response is not valid JSON.} \\
\texttt{Output ONLY:} \\
\texttt{One JSON with confidence\_score, current\_mode, input\_rationale, and hypothesis keys ...} \\
\texttt{OR} \\
\texttt{SOLVE: ```python ... ```} \\
            }
        }
    }
     \\ \\

    \theutterance \stepcounter{utterance}  
    & \multicolumn{4}{p{0.6\linewidth}}{
        \cellcolor[rgb]{0.561,0.737,0.561}{
            \makecell[{{p{\linewidth}}}]{
                \texttt{\tiny{[P1$\rangle$GM]}}
                \texttt{```json} \\
\texttt{\{} \\
\texttt{  "current\_mode": "probing",} \\
\texttt{  "input\_rationale": "Given the contradictions in previous attempts, I'll try a simpler hypothesis: the function returns the count of elements that are greater than the previous element (i.e., the number of 'increases' in the sequence). This fits the preloaded example ({[}5, {-}5, 5{]}, 2) if we count the transitions from {-}5 to 5 (1 increase) and ignore the first element, but this doesn't match. Alternatively, it might count the number of times the sequence changes direction. Given the time constraints, I'll submit a SOLVE attempt based on the preloaded example and the latest hint ({[}0{]}, 0).",} \\
\texttt{  "hypothesis": "unknown",} \\
\texttt{  "confidence\_score": 30} \\
\texttt{\}} \\
\texttt{```} \\
\\ 
\texttt{SOLVE:} \\
\texttt{```python} \\
\texttt{def f(items):} \\
\texttt{    """Returns the number of local maxima in the list, where a local maximum is an element greater than its neighbors."""} \\
\texttt{    if len(items) == 1:} \\
\texttt{        return 0} \\
\texttt{    count = 0} \\
\texttt{    for i in range(len(items)):} \\
\texttt{        if i == 0:} \\
\texttt{            if items{[}i{]} > items{[}i+1{]}:} \\
\texttt{                count += 1} \\
\texttt{        elif i == len(items) {-} 1:} \\
\texttt{            if items{[}i{]} > items{[}i{-}1{]}:} \\
\texttt{                count += 1} \\
\texttt{        else:} \\
\texttt{            if items{[}i{]} > items{[}i} \\
            }
        }
    }
    & & \\ \\

    \theutterance \stepcounter{utterance}  
    & & & \multicolumn{4}{p{0.6\linewidth}}{
        \cellcolor[rgb]{0.9,0.9,0.9}{
            \makecell[{{p{\linewidth}}}]{
                \texttt{\tiny{[P1$\langle$GM]}}
                \texttt{Game Violation: ParseError: Response is not valid JSON.} \\
\texttt{Output ONLY:} \\
\texttt{One JSON with confidence\_score, current\_mode, input\_rationale, and hypothesis keys ...} \\
\texttt{OR} \\
\texttt{SOLVE: ```python ... ```} \\
            }
        }
    }
     \\ \\

    \theutterance \stepcounter{utterance}  
    & \multicolumn{4}{p{0.6\linewidth}}{
        \cellcolor[rgb]{0.561,0.737,0.561}{
            \makecell[{{p{\linewidth}}}]{
                \texttt{\tiny{[P1$\rangle$GM]}}
                \texttt{```json} \\
\texttt{\{} \\
\texttt{  "current\_mode": "confirming",} \\
\texttt{  "input\_rationale": "After analyzing all examples, the function appears to count the number of times the sequence changes from increasing to decreasing or vice versa (i.e., the number of 'peaks' and 'valleys'). The preloaded example ({[}5, {-}5, 5{]}, 2) fits because there are two direction changes (5→{-}5 is decreasing, {-}5→5 is increasing). The example ({[}1, 2, 3{]}, {-}1) suggests it returns {-}1 for strictly increasing/decreasing sequences. The example ({[}0{]}, 0) fits as a single element has no changes. I will now submit the correct implementation.",} \\
\texttt{  "hypothesis": "def f(items):\textbackslash{}n    \textbackslash{}"\textbackslash{}"\textbackslash{}"Returns the number of direction changes in the sequence, or {-}1 if strictly increasing/decreasing.\textbackslash{}"\textbackslash{}"\textbackslash{}"\textbackslash{}n    if len(items) <= 1:\textbackslash{}n        return 0\textbackslash{}n    direction\_changes = 0\textbackslash{}n    prev\_diff = 0\textbackslash{}n    for i in range(1, len(items)):\textbackslash{}n        curr\_diff = items{[}i{]} {-} items{[}i{-}1{]}\textbackslash{}n        if curr\_diff != 0:\textbackslash{}n            if prev\_diff * curr\_diff < 0:\textbackslash{}n                direction\_changes += 1\textbackslash{}n            prev\_diff = curr\_diff\textbackslash{}n    if direction\_changes == 0 and prev\_diff != 0:\textbackslash{}n} \\
            }
        }
    }
    & & \\ \\

    \theutterance \stepcounter{utterance}  
    & & & \multicolumn{4}{p{0.6\linewidth}}{
        \cellcolor[rgb]{0.9,0.9,0.9}{
            \makecell[{{p{\linewidth}}}]{
                \texttt{\tiny{[P1$\langle$GM]}}
                \texttt{Game Violation: ParseError: Response is not valid JSON.} \\
\texttt{Output ONLY:} \\
\texttt{One JSON with confidence\_score, current\_mode, input\_rationale, and hypothesis keys ...} \\
\texttt{OR} \\
\texttt{SOLVE: ```python ... ```} \\
            }
        }
    }
     \\ \\

    \theutterance \stepcounter{utterance}  
    & \multicolumn{4}{p{0.6\linewidth}}{
        \cellcolor[rgb]{0.561,0.737,0.561}{
            \makecell[{{p{\linewidth}}}]{
                \texttt{\tiny{[P1$\rangle$GM]}}
                \texttt{SOLVE:} \\
\texttt{```python} \\
\texttt{def f(items):} \\
\texttt{    """Returns the number of times the sequence changes direction (peaks + valleys), or {-}1 if strictly increasing/decreasing."""} \\
\texttt{    if len(items) <= 1:} \\
\texttt{        return 0} \\
\texttt{    direction\_changes = 0} \\
\texttt{    prev\_diff = items{[}1{]} {-} items{[}0{]}} \\
\texttt{    for i in range(2, len(items)):} \\
\texttt{        curr\_diff = items{[}i{]} {-} items{[}i{-}1{]}} \\
\texttt{        if curr\_diff * prev\_diff < 0:} \\
\texttt{            direction\_changes += 1} \\
\texttt{        if curr\_diff != 0:} \\
\texttt{            prev\_diff = curr\_diff} \\
\texttt{    if direction\_changes == 0 and abs(prev\_diff) > 0:} \\
\texttt{        return {-}1} \\
\texttt{    return direction\_changes} \\
\texttt{```} \\
            }
        }
    }
    & & \\ \\

    \theutterance \stepcounter{utterance}  
    & & & \multicolumn{2}{p{0.3\linewidth}}{
        \cellcolor[rgb]{0.9,0.9,0.9}{
            \makecell[{{p{\linewidth}}}]{
                \texttt{\tiny{[GM$|$GM]}}
                \texttt{Game Verdict: LOSS} \\
            }
        }
    }
    & & \\ \\

    \theutterance \stepcounter{utterance}  
    & & & \multicolumn{2}{p{0.3\linewidth}}{
        \cellcolor[rgb]{0.9,0.9,0.9}{
            \makecell[{{p{\linewidth}}}]{
                \texttt{\tiny{[GM$|$GM]}}
                \texttt{Incorrect.} \\
\texttt{Failed Cases:} \\
\texttt{Input: {[}{[}5, {-}5, 5{]}{]} | Expected: 2 | Got: 1} \\
\texttt{Input: {[}{[}1, 2, 3{]}{]} | Expected: 0 | Got: {-}1} \\
\texttt{Input: {[}{[}3, 2, 1{]}{]} | Expected: 0 | Got: {-}1} \\
            }
        }
    }
    & & \\ \\

    \theutterance \stepcounter{utterance}  
    & & & \multicolumn{2}{p{0.3\linewidth}}{
        \cellcolor[rgb]{0.9,0.9,0.9}{
            \makecell[{{p{\linewidth}}}]{
                \texttt{\tiny{[GM$|$GM]}}
                \texttt{def count\_adjacent\_opposite\_pairs(items: List{[}int{]}) {-}> int:} \\
\texttt{    return sum(1 for i in range(len(items) {-} 1) if items{[}i{]} == {-}items{[}i + 1{]})} \\
            }
        }
    }
    & & \\ \\

    \theutterance \stepcounter{utterance}  
    & & & \multicolumn{2}{p{0.3\linewidth}}{
        \cellcolor[rgb]{0.9,0.9,0.9}{
            \makecell[{{p{\linewidth}}}]{
                \texttt{\tiny{[GM$|$GM]}}
                \texttt{def f(items):} \\
\texttt{    """Returns the number of times the sequence changes direction (peaks + valleys), or {-}1 if strictly increasing/decreasing."""} \\
\texttt{    if len(items) <= 1:} \\
\texttt{        return 0} \\
\texttt{    direction\_changes = 0} \\
\texttt{    prev\_diff = items{[}1{]} {-} items{[}0{]}} \\
\texttt{    for i in range(2, len(items)):} \\
\texttt{        curr\_diff = items{[}i{]} {-} items{[}i{-}1{]}} \\
\texttt{        if curr\_diff * prev\_diff < 0:} \\
\texttt{            direction\_changes += 1} \\
\texttt{        if curr\_diff != 0:} \\
\texttt{            prev\_diff = curr\_diff} \\
\texttt{    if direction\_changes == 0 and abs(prev\_diff) > 0:} \\
\texttt{        return {-}1} \\
\texttt{    return direction\_changes} \\
            }
        }
    }
    & & \\ \\

    \theutterance \stepcounter{utterance}  
    & & & \multicolumn{2}{p{0.3\linewidth}}{
        \cellcolor[rgb]{0.9,0.9,0.9}{
            \makecell[{{p{\linewidth}}}]{
                \texttt{\tiny{[GM$|$GM]}}
                \texttt{Mode: passive\_labeled\_pairs} \\
\texttt{Accuracy: 0.170} \\
\texttt{Information requests: 2/15} \\
\texttt{Revealed passive examples: 3} \\
\texttt{Efficiency: 0.933} \\
\texttt{Slack: 13} \\
\texttt{Observed{-}pair consistency: 0.500} \\
\texttt{Membership balance: 2 true / 1 false} \\
\texttt{Parse errors: 5} \\
\texttt{Runtime errors: 0} \\
            }
        }
    }
    & & \\ \\

\end{supertabular}
}

\end{document}
